% !TeX program = xelatex
\documentclass[11pt]{article}

\usepackage[a4paper, margin=1in]{geometry}
\usepackage{amsmath, amssymb, amsthm}
\usepackage{stmaryrd}
\usepackage{listings}
\usepackage{xcolor}
\usepackage{hyperref}
\usepackage{microtype}
\usepackage{fontspec}
\usepackage{booktabs}
\setmonofont{DejaVu Sans Mono}[Scale=0.85]

\sloppy
\setlength{\emergencystretch}{3em}

\newcommand{\btt}[1]{\texttt{\detokenize{#1}}}

% ---- Lean code listing style -------------------------------------------
\lstdefinelanguage{lean}{
  morekeywords={theorem, lemma, def, by, fun, let, in, do, match, with,
    if, then, else, return, import, open, namespace, end, structure,
    inductive, partial, private, noncomputable, syntax, elab, elab_rules,
    macro, macro_rules, tactic, conv, command, example, have, show, intro,
    apply, simp, rw, ring, all_goals, repeat, decide, native_decide,
    fin_cases, norm_num, exact, set_option},
  sensitive=true,
  morecomment=[l]{--},
  morecomment=[s]{/-}{-/},
  morestring=[b]"
}
\lstset{
  inputencoding=utf8,
  language=lean,
  basicstyle=\ttfamily\footnotesize,
  keywordstyle=\color{blue!70!black}\bfseries,
  commentstyle=\color{gray}\itshape,
  stringstyle=\color{red!60!black},
  showstringspaces=false,
  breaklines=true,
  columns=fullflexible,
  frame=single,
  framesep=4pt,
  rulecolor=\color{black!30},
  backgroundcolor=\color{gray!6},
  xleftmargin=6pt,
  xrightmargin=6pt,
  aboveskip=0.8\baselineskip,
  belowskip=0.8\baselineskip
}
% -----------------------------------------------------------------------

\title{From Flagmatic Certificate to Verified Lean Proof}
\author{}
\date{}

\begin{document}
\maketitle

\section{Introduction}

The flag algebra method, introduced by Razborov, has become a workhorse of
extremal combinatorics. A typical proof of a density inequality
$d(F;G) \le \alpha$ in some forbidden-subgraph class proceeds by exhibiting
a positive semidefinite matrix $M$ and a vector $v$ of subflags such that
the quadratic form $v^\top M v$ closes a gap in a linear combination of
flag densities. Software such as flagmatic produces these certificates
automatically: it solves an SDP, rounds the floating-point solution to
exact rationals, and outputs a JSON file containing the matrices, the
chosen subflag types, and a manifest of the admissible host flags.

What is normally \emph{not} in scope for such tools is a machine-checked
proof. A flagmatic certificate is a body of numerical data; to turn it
into a verified theorem one must transcribe every number into a proof
assistant, reproduce the bookkeeping that links subflag indices to host
flags, and check positive semidefiniteness, density tables, and product
expansions inside the kernel.

This paper documents \emph{a framework that closes this gap}. Given a
flagmatic-style certificate JSON and a target file path, one command
\begin{lstlisting}[language={}]
python flagmatic_to_lean.py gen-skeleton <cert>.json <target>.lean \
       --namespace <Name>
\end{lstlisting}
produces a complete Lean~4 file that the kernel accepts. The output
contains the SDP matrices, an exact-rational $\mathrm{LDL}^\top$
positive-semidefiniteness proof for each, the bulk-load commands for the
relevant flag/density/product tables, the subflag vector declarations,
and the body of the main theorem, all assembled so that
\texttt{lake build} succeeds with zero \texttt{sorry}.

The framework rests on three layers built together: a layer of generated
flag definitions and density tables (Section~\ref{sec:flag-defs},
Section~\ref{sec:density-tables}), a layer of proof-time tactics that
consume them (Section~\ref{sec:tactics}), and the certificate-to-Lean
translation script that is the focus of this paper
(Sections~\ref{sec:format}--\ref{sec:body}). Section~\ref{sec:background}
introduces the infrastructure that the rest of the paper builds on, so
this article is self-contained.

We use the names from the public layout of the repository throughout. The
Lean side lives under \texttt{LeanFlagAlgebras/}; the translation script is
\texttt{LeanFlagAlgebras/Flagmatic/flagmatic\_to\_lean.py}; certificates
sit in \texttt{LeanFlagAlgebras/Flagmatic/Certificates/}, and generated
proofs in \texttt{LeanFlagAlgebras/Flagmatic/}.

\section{Background: the Lean framework underneath}
\label{sec:background}

This section introduces the pieces of the Lean side of the framework that
the rest of the paper builds on: the generated flag definitions, the
pre-computed density tables, the proof-time tactics, and the
four-part proof shape that ties them together.

\subsection{Generated flag definitions}
\label{sec:flag-defs}

Every flag in the development carries its identifying indices in its name.
For an unlabeled flag (``empty-typed'', in the framework's vocabulary) of
size $n$ at canonical index $i$ we have
\begin{itemize}
  \item \texttt{Flag\_n\_0\_0\_i} -- the flag as an object in the flag
    category, and
  \item \texttt{FlagAlgebra\_n\_0\_0\_i} -- the corresponding generator of
    the flag algebra, defined as $\llbracket
    \mathrm{unitVector}\,\langle n, \texttt{Flag\_n\_0\_0\_i}\rangle\rrbracket$.
\end{itemize}
For a flag of size $n$ on a $k$-vertex type of canonical type-index $m$,
similarly,
\begin{itemize}
  \item \texttt{FlagType\_k\_m} -- the type as a labeled graph,
  \item \texttt{Flag\_n\_k\_m\_i}, \texttt{FlagAlgebra\_n\_k\_m\_i} --
    flag and algebra generator on that type.
\end{itemize}

These names are not written by hand. Two meta-commands, defined in
\texttt{Flags/FlagLoader.lean}, register them in bulk from JSON files:
\begin{itemize}
  \item \texttt{load\_empty\_typed\_flags "graphs\_n.json"} reads the
    canonical enumeration of all unlabeled graphs on $n$ vertices and
    registers \texttt{Flag\_n\_0\_0\_i}, \texttt{FlagAlgebra\_n\_0\_0\_i},
    together with the bookkeeping lemmas
    \texttt{flagSet\_n\_0\_0\_val\_eq} (the list of all such flags) and
    \texttt{flagSet\_n\_0\_0\_eq\_univ}. Each lemma is closed by
    \texttt{native\_decide}, so the enumeration is re-verified inside the
    kernel.
  \item \texttt{load\_flags "flags\_n\_k\_m.json"} is the labeled
    counterpart, registering \texttt{Flag\_n\_k\_m\_i},
    \texttt{FlagAlgebra\_n\_k\_m\_i}, and the analogous
    \texttt{flagSet\_n\_k\_m\_*} lemmas, plus two \texttt{@[simp]} lemmas
    (\texttt{unlabel\_n\_k\_m\_i}, \texttt{downward\_n\_k\_m\_i}) that
    relate a labeled flag to the underlying unlabeled flag.
\end{itemize}
The input JSON files come from two small Python scripts:
\texttt{Flags/generate\_graphs.py} (for the unlabeled case, using
NetworkX's graph atlas with a lexicographic-minimal canonicalization
pass) and \texttt{Flags/generate\_flags.py} (for the labeled case, which
additionally records each flag's type embedding and downward normalizing
factor). The repository ships precomputed JSON for sizes up to
$n=5$, loaded together in \texttt{Flags/FlagDef.lean}, so identifiers
such as \texttt{Flag\_5\_0\_0\_17} are available project-wide.

\subsection{Pre-computed density tables}
\label{sec:density-tables}

Three further meta-commands, defined in \texttt{Flags/Densities/}, register
the numerical content needed by forbid-conditioned proofs.

\paragraph{\texttt{load\_forbid\_density\_theorems "<file>.json"}}
For a fixed forbidden subgraph $H$, this command reads a JSON file that
lists the indices of size-$n$ empty-typed flags that are $H$-free, and
registers, for every flag $F_i$, an \texttt{@[simp]} lemma stating
$\mathrm{flagDensity}_1\,H.\mathrm{toFinFlag}\,F_i = 0$ if $F_i$ is
$H$-free, and $\ne 0$ otherwise. The supported tags currently are $K_3$,
$K_4$, and $K_5$; each branch is closed by
\texttt{flagDensity\textsubscript{1}\_eq\_sym2EmptyTypeFlagDensity\textsubscript{1}}
followed by \texttt{native\_decide}.

\paragraph{\texttt{load\_flag\_pair\_density\_theorems "<file>.json"}}
Reads a JSON file produced by \texttt{calculate\_densities.py}
containing pair densities $\mathrm{flagDensity}_2(F_{p_1}, F_{p_2}; F_h)$
of two pattern flags inside a host flag. For each row it registers an
\texttt{@[simp]} lemma stating the exact rational value of that pair
density, closed by \texttt{native\_decide}.

\paragraph{\texttt{load\_forbid\_mul\_theorems "<file>.json"}}
Consumes the \emph{same} JSON as the previous command and produces, for
every pair $(i,j)$ of pattern indices with $i \le j$, an expansion
theorem
\[
  F_i \cdot F_j \;=_{H\text{-forbid}}\; \sum_h c_{ijh}\,F_h,
\]
named \texttt{flagMul\_FlagAlgebra\_<patTag>\_i\_FlagAlgebra\_<patTag>\_j}.
The proof rewrites with \texttt{unitVector\_quot\_mul\_forbidEq\_sum},
expands the host sum with the \texttt{flagSet\_*\_val\_eq} lemmas, and
drops the forbidden host flags by \texttt{simp}.

The JSON inputs themselves come from two more Python scripts:
\texttt{Flags/Densities/calculate\_densities.py} computes the pair-density
tables exhaustively over rational arithmetic, and
\texttt{Flags/Densities/gen\_free\_indices.py} computes the $H$-free index
lists. As with the flag enumeration, all numerical content is re-verified
in Lean; the scripts are convenience, not trusted-base.

\subsection{Proof tactics}
\label{sec:tactics}

The names registered by Sections~\ref{sec:flag-defs} and
\ref{sec:density-tables} are picked up by a small set of metaprogrammed
tactics. Each tactic looks at the goal's expression, synthesizes the
required \texttt{flagSet\_*}, \texttt{downward\_*}, or
\texttt{flagMul\_*} name from the indices it sees, and applies it.

\paragraph{\texttt{expand\_one\_at n}}
(\texttt{API/Basic.lean}.) Under a forbid condition, rewrites the unit
$1 \in \mathrm{FlagAlgebra}$ as the weighted sum of empty-typed flags of
size $n$ that are $H$-free. The same call works at every size; the tactic
synthesizes \texttt{flagSet\_n\_0\_0\_eq\_univ} and
\texttt{flagSet\_n\_0\_0\_val\_eq} from its argument.

\paragraph{\texttt{reduce\_downward\_flagmul}}
(\texttt{API/ReduceFlagMul.lean}.) Handles \texttt{forbidLE} goals whose
left-hand side is a sum of terms of one of three shapes:
\texttt{downward(}$c \cdot (A \cdot B)$\texttt{)},
\texttt{downward(}$A \cdot B$\texttt{)} (with the scalar absorbed), or a
plain flag term. For each leading product the tactic looks up the
matching \texttt{flagMul\_<A>\_<B>} theorem (Section~\ref{sec:density-tables}),
substitutes it via \texttt{downward\_forbidEq\_equal\_flags}, and moves
the resulting linear combination of base flags to the right-hand side.
Plain flag terms are moved directly. The loop is bounded by a fuel of
$256$ steps and is preceded by a normalizing
\texttt{simp only [downward\_add, add\_assoc]} pass.

\paragraph{\texttt{prove\_flag\_expand n}}
(\texttt{Utils/FlagExpansionTactic.lean}.) Closes a goal of the form
$F = \sum_j c_j\,F_j$ where the right-hand side is the size-$n$ expansion
of $F$ as a sum of empty-typed flags. The proof script unfolds $F$ via
\texttt{unitVector\_eqv\_densityFlagSum}, expands the resulting sum using
the \texttt{flagSet\_n\_*} lemmas, and closes by \texttt{ring\_nf}. The
last step relies on the \texttt{flagDensity}$_1$ values being already
present as \texttt{@[simp]} lemmas; the translation script (Section~8)
emits them automatically when this tactic is needed.

\paragraph{\texttt{flag\_nonneg}}
(\texttt{API/Basic.lean}.) Closes a goal of the form
$0 \le c_1\,\llbracket\mathrm{unitVector}\,F_1\rrbracket + \cdots +
c_k\,\llbracket\mathrm{unitVector}\,F_k\rrbracket$ with $c_i \ge 0$. The
macro introduces an arbitrary positive homomorphism, distributes it over
the sum, splits via \texttt{repeat apply add\_nonneg}, and discharges each
leaf with \texttt{positiveHom\_unitVector\_ge\_zero}.

\paragraph{\texttt{ac\_sort\_rhs\_pipeline}}
(\texttt{Utils/SortTactic.lean}.) Normalizes the right-hand side of an
equation between linear combinations of flags by stable-sorting the
summands on their base-flag index, merging coefficients of the same base
flag, and discharging the resulting equation up to add-AC.

\subsection{The proof shape}
\label{sec:proof-shape}

Combining the elements above, every main API proof in the framework
follows the same four-part shape.
\begin{enumerate}
  \item \textbf{Part 1.} Per-block SDP certificate matrices and their
    $\mathrm{LDL}^\top$-based positive-semidefiniteness proofs.
  \item \textbf{Part 2.} The bulk-load meta-commands of
    Section~\ref{sec:density-tables}, registering forbid-free indices,
    pair densities, and product expansions for the relevant
    (host, pattern) pairs.
  \item \textbf{Part 3.} The subflag type and the flag vectors $v_t$ that
    each matrix $M_t$ is paired with.
  \item \textbf{Part 4.} The theorem statement and proof body.
\end{enumerate}
This shape is regular enough that the translation script we describe in
the rest of the paper can produce a fully populated instance of it from
a flagmatic certificate alone.

\section{The certificate format}
\label{sec:format}

The input to the framework is a JSON file in the ``sparse SDP output''
format produced by flagmatic. The schema is fixed: the
following twelve fields, exactly. We use Mantel's theorem as a running
example.

\begin{lstlisting}[language={}]
{
  "description":               "2-graph; maximize 2:12 density; forbid 3:121323",
  "bound":                     "1/2",
  "order_of_admissible_graphs": 3,
  "number_of_admissible_graphs": 3,
  "admissible_graphs":         ["3:", "3:12", "3:1213"],
  "admissible_graph_densities": [0, "1/3", "2/3"],
  "number_of_types":           1,
  "types":                     ["1:"],
  "numbers_of_flags":          [2],
  "flags":                     [["2:(1)", "2:12(1)"]],
  "qdash_matrices":            [[[2]]],
  "r_matrices":                [[["1/2"], ["-1/2"]]]
}
\end{lstlisting}

\paragraph{Flagmatic strings.}
Graphs are encoded as compact strings of the form
\texttt{N:e\textsubscript{1}e\textsubscript{2}\ldots}\,, where $N$ is the
vertex count and each $e_i$ is a two-digit pair of vertices
(1-indexed, single-digit) forming one edge. Thus
\texttt{"3:1213"} is the path $1$\,-$2$, $1$\,-$3$. A sigma-flag uses
the form \texttt{m:edges(k)}, where the parenthesized integer $k$ is the
type size; by convention the first $k$ vertices are the type vertices in
label order. Hence \texttt{"3:12(2)"} is a $3$-vertex flag on a type
of size $2$ (vertices $1,2$) with an additional edge between the type
vertices.

\paragraph{The matrix data.}
The certificate carries each SDP block as a pair
$(Q_t', R_t)$, not as the assembled positive semidefinite matrix.
\texttt{r\_matrices[t]} is the (number of flags $\times$ rank-of-block)
matrix $R_t$, stored row by row; \texttt{qdash\_matrices[t]} is the
inner positive semidefinite matrix $Q_t'$, stored in
\emph{upper-triangular row form}, that is, row $i$ contains
$Q_t'[i,i], Q_t'[i,i+1], \ldots, Q_t'[i,n-1]$. The assembled matrix is
\[
  M_t \;=\; R_t \cdot Q_t' \cdot R_t^\top,
\]
of dimension (number of flags) $\times$ (number of flags). For Mantel,
$R_1 = (1/2, -1/2)^\top$ and $Q_1' = [\,2\,]$ produce
$M_1 = \begin{pmatrix} 1/2 & -1/2 \\ -1/2 & 1/2 \end{pmatrix}$.

\paragraph{The admissible-graph manifest.}
\texttt{admissible\_graphs} lists every $N$-vertex graph that is
$H$-free, in flagmatic's canonical order. \texttt{admissible\_graph\_-
densities} gives, in the same order, the density of the objective
graph (the one named in the \texttt{description} field's
\texttt{maximize} clause) in each admissible host. This manifest carries
the information the SDP solver actually used and is the reference both
sides of the eventual inequality refer to.

\paragraph{The description field.}
The \texttt{description} string has a fixed grammar:
\texttt{2-graph; maximize <objective\_str> density; forbid <forbid\_str>}.
The translation script parses both names out of it, both for identifier
lookup (Section~5) and for displaying a docstring on the generated
theorem.

\section{The cert-to-Lean pipeline}

The script \texttt{flagmatic\_to\_lean.py} exposes five subcommands:

\begin{lstlisting}[language={}]
python flagmatic_to_lean.py inspect       <cert>.json
python flagmatic_to_lean.py check-deps    <cert>.json
python flagmatic_to_lean.py gen-skeleton  <cert>.json <target>.lean [opts]
python flagmatic_to_lean.py gen-matrices  <cert>.json <target>.lean
python flagmatic_to_lean.py gen-vectors   <cert>.json <target>.lean
\end{lstlisting}

\paragraph{\texttt{inspect}.} A read-only diagnostic. Prints, for every
flagmatic string in the certificate (admissible graphs, types,
$\sigma$-flags), the Lean identifier it resolves to via the mapping of
Section~5. Useful as a sanity check on a fresh certificate.

\paragraph{\texttt{check-deps}.} Reports which JSON files
(Sections~\ref{sec:flag-defs}, \ref{sec:density-tables}) the certificate
will reference, marks each as present or missing on disk, and emits the
exact \texttt{import} and \texttt{load\_*} lines that the generated Lean
file will use. Exits with code $1$ if any required file is missing,
making it convenient as a CI guard.

\paragraph{\texttt{gen-skeleton}.} The main entry point. Produces a
complete Lean file containing the four parts of
Section~\ref{sec:proof-shape}: imports and opens, namespace, load
commands, per-block matrices and PSD proofs, $\sigma$ and $v$ vector
declarations, and the body of the main theorem. The output is a
self-contained module that \texttt{lake build} accepts.

\paragraph{\texttt{gen-matrices}, \texttt{gen-vectors}.} Append-only
helpers that emit just one part of the four; useful when an existing
file is being extended rather than freshly generated. Both are
single-block-aware and respect the same identifier conventions as
\texttt{gen-skeleton}.

The standard workflow is
\begin{lstlisting}[language={}]
inspect -> check-deps -> gen-skeleton -> lake build
\end{lstlisting}
which catches identifier mismatches and missing dependencies before any
Lean compilation is attempted.

\section{Resolving flagmatic strings to Lean identifiers}
\label{sec:mapping}

The certificate refers to graphs by flagmatic strings; the Lean side
refers to them by indexed names registered with the meta-commands of
Section~\ref{sec:flag-defs}. The translation script's first job is to
match these.

\paragraph{Parsing.}
Each flagmatic string is parsed into a triple
$(n, E, \mathit{labels})$, where $E$ is a frozen set of $0$-indexed
edges and $\mathit{labels}$ is either \texttt{None} (for an unlabeled
graph or a type) or a tuple of label positions (for a sigma-flag). The
parser rejects odd-length edge fields, self-loops, duplicate edges, and
out-of-range vertex indices.

\paragraph{Unlabeled lookup.}
For an unlabeled flagmatic string with size $n$ and edge set $E$, the
translation script loads the canonical list of graphs from
\texttt{graphs\_n.json} and searches for an index $i$ such that the
$i$-th canonical graph $G_i$ is isomorphic to $(n, E)$. The search
brute-forces over all $n!$ vertex permutations -- fine in practice since
$n \le 5$ in the current development. The result is the index $i$ that
appears in the Lean names \texttt{Flag\_n\_0\_0\_i} and
\texttt{FlagAlgebra\_n\_0\_0\_i}.

\paragraph{Sigma-flag lookup.}
For a sigma-flag \texttt{m:edges(k)} we first identify the underlying
unlabeled graph as above, then load
\texttt{flags\_m\_k\_<typeIdx>.json}, which records, for every
$\sigma$-flag of that shape, the canonical positions
\texttt{type\_indices} of the type vertices. The script then searches
for a permutation that simultaneously sends the input edges to the
stored edges and the input type positions (the first $k$ vertices, by
convention) to the stored \texttt{type\_indices}. The result is the
index $i$ in \texttt{FlagAlgebra\_m\_k\_typeIdx\_i}.

\paragraph{Forbid resolution.}
The \texttt{description} field of the certificate names the forbidden
graph by a flagmatic string, but the proof body needs a concrete Lean
term such as \texttt{K3.toFinFlag}. To map between the two, the script
parses \texttt{Forbid/CommonGraphs.lean} with the regular expression
\begin{lstlisting}[language={}]
lemma <Name>_toFinFlag_eq : <Name>.toFinFlag = ⟨n, Flag_n_0_0_i⟩
\end{lstlisting}
For each match it records \texttt{(<Name>, n, i)}, then resolves $i$
against \texttt{graphs\_n.json} to obtain the canonical edge set of the
named forbid graph. The certificate's forbid string is then
isomorphism-matched against this table. The match is canonical (vertex
permutations exhausted as above), so e.g.\ a triangle described as
\texttt{3:121323}, \texttt{3:122331}, or any other reordering all resolve
to \texttt{K3}. Adding a new forbid graph -- the topic of Section~11 --
amounts to extending \texttt{CommonGraphs.lean} with two lines and
re-running the data scripts; no change to the Python is needed.

\section{Matrix assembly and PSD proofs (generates Part 1)}
\label{sec:matrices}

For each SDP block $t$, the translation script produces a six-lemma cluster
in Lean. The shape, taken from the generated \texttt{K3forbidC4.lean}
(K3-free, $N=4$, two blocks), is

\begin{lstlisting}
def M₁ : Matrix (Fin 4) (Fin 4) ℚ :=
  !![(3 / 8 : ℚ), 0, 0, (-3 / 8 : ℚ);
     0, 0, 0, 0;
     0, 0, 0, 0;
     (-3 / 8 : ℚ), 0, 0, (3 / 8 : ℚ)]
noncomputable def M₁_real : Matrix (Fin 4) (Fin 4) ℝ :=
  ratMatrixToReal M₁

def dM₁ : Fin 4 → ℚ := ![(3 / 8 : ℚ), 0, 0, 0]
def LM₁ : Matrix (Fin 4) (Fin 4) ℚ :=
  !![1, 0, 0, 0;
     0, 1, 0, 0;
     0, 0, 1, 0;
     (-1 : ℚ), 0, 0, 1]

lemma dM₁_nonneg (i : Fin 4) : 0 ≤ dM₁ i := by
  fin_cases i <;> norm_num [dM₁]
lemma M₁_eq_LDL : M₁ = LM₁ * Matrix.diagonal dM₁ * LM₁ᵀ := by
  decide +kernel
theorem M₁_posSemidef : M₁.PosSemidef :=
  posSemidef_of_eq_mul_diagonal_mul_transpose dM₁_nonneg M₁_eq_LDL
-- The same decomposition is then lifted over ℝ to obtain M₁_real_posSemidef.
\end{lstlisting}

\paragraph{Assembly.}
The script first reconstructs $M_t = R_t \cdot Q_t' \cdot R_t^\top$
exactly. Both $R_t$ and $Q_t'$ are read out of the certificate as Python
\texttt{Fraction}s, with $Q_t'$ being unfolded from upper-triangular row
form into a full symmetric matrix.

\paragraph{$\mathrm{LDL}^\top$ decomposition.}
The decomposition is computed by a textbook elimination loop over
\texttt{Fraction}s, with a small modification to handle the rank
deficiency that the assembled matrix typically has. Concretely, for each
column index $i$ in turn:
\[
  d_i \;=\; M_t[i,i] \;-\; \sum_{k<i} L[i,k]^2 \, d_k,
  \qquad
  L[j,i] \;=\;
  \frac{M_t[j,i] \;-\; \sum_{k<i} L[j,k]\,L[i,k]\,d_k}{d_i}
  \quad (j>i).
\]
If $d_i = 0$ at some column then the numerator on the right must also be
zero for every $j > i$; if not, the input matrix is rejected as not
positive semidefinite. When the numerator vanishes the entry $L[j,i]$ is
set to $0$.

\paragraph{Lean re-verification.}
None of this computation is in the trusted base. The Lean kernel
re-verifies the equality $M_t = L \cdot \mathrm{diag}(d) \cdot L^\top$
on its own with \texttt{decide +kernel}, and turns it into PSD-ness via
a generic lemma. A miscomputation in the script -- a wrong $L$ or $d$
-- surfaces as a compile error on \texttt{M\_eq\_LDL}, not as a silently
accepted bad PSD claim.

\paragraph{Lifting to $\mathbb{R}$.}
The proof inside the framework's proof body needs $M_{t,\mathbb{R}}$ to
be PSD over the reals, not just over the rationals. A second cluster
lifts the rational decomposition through \texttt{ratMatrixToReal} and a
real-coercion-aware generic lemma; the Lean text is mechanical and is
emitted verbatim by the script.

\section{Generating the load commands (generates Part 2)}

For each block $t$ the certificate fixes the host size $N$, the type
$\sigma_t$ (of size $k$ and canonical index $\tau_t$), and the
$\sigma$-flag dimension $m_t$. Together with the forbid tag $H$ these
determine exactly which JSON files the proof needs:
\begin{itemize}
  \item \texttt{graphs\_N\_<H>\_free\_indices.json}, registered once by
    \texttt{load\_forbid\_density\_theorems};
  \item for each block, \texttt{density\_N\_k\_$\tau_t$\_from\_$m_t$\_k\_$\tau_t$\_forbid\_<H>.json},
    registered by both
    \texttt{load\_flag\_pair\_density\_theorems} and
    \texttt{load\_forbid\_mul\_theorems} -- the same JSON is the source
    of both density values and product expansions.
\end{itemize}
For the running Mantel example ($N=3$, one block on the empty type of
size $1$, $m_1 = 2$, $H = K_3$) the script emits
\begin{lstlisting}
load_forbid_density_theorems
  "LeanFlagAlgebras/Flags/Densities/graphs_3_K3_free_indices.json"
load_flag_pair_density_theorems
  "LeanFlagAlgebras/Flags/Densities/density_3_1_0_from_2_1_0_forbid_K3.json"
load_forbid_mul_theorems
  "LeanFlagAlgebras/Flags/Densities/density_3_1_0_from_2_1_0_forbid_K3.json"
\end{lstlisting}

\paragraph{Strict dependency checking.}
An earlier version of the code accepted, as a fallback, density files
without the \texttt{\_forbid\_<H>} suffix. That fallback was unsafe:
two computations under different forbid graphs both default to the same
file name, so one of them silently overwrites the other and the
overwriting set of densities is silently loaded by the proof. We
removed the fallback. The translation script now demands an
explicit \texttt{\_forbid\_<H>} suffix (or \texttt{\_no\_forbid} for the
unconditioned case) for every density file, and \texttt{check-deps}
reports any missing variant rather than silently substituting
another.

\section{The main theorem: branches A and B (generates Part 4)}
\label{sec:body}

The body of the main theorem follows one of two templates, selected by
comparing the objective flag's vertex count $n_{\mathit{obj}}$ to the
host size $N$.

\subsection{Branch A: $n_{\mathit{obj}} = N$}

When the objective is already at host size -- as in $C_4$-Tur\'an
($n_{\mathit{obj}} = 4 = N$) or Erd\H{o}s pentagon
($n_{\mathit{obj}} = 5 = N$) -- the body is a fixed eight-step sequence,
parametrized only by $N$, the block count $T$, and the per-block
dimensions $\{n_t\}_{t=1}^{T}$. The output for the K3forbidC4 example
($N=4$, $T=2$, $n_1=4$, $n_2=3$) is

\begin{lstlisting}
theorem k3forbidC4_flagAlgebra
    : FlagAlgebra_4_0_0_8 ≤[K3.toFinFlag] (3 / 8 : ℝ) • (1 : FlagAlgebra ∅ₜ) := by
  have quadraticForm_trans : FlagAlgebra_4_0_0_8 ≤[K3.toFinFlag]
        FlagAlgebra_4_0_0_8
        + ⟦flagQuadraticForm M₁_real v₁⟧₀
        + ⟦flagQuadraticForm M₂_real v₂⟧₀ := by
    apply forbidLE_add_QuadraticForm M₂_real M₂_real_posSemidef v₂
    apply forbidLE_add_QuadraticForm M₁_real M₁_real_posSemidef v₁
    exact forbidLE_refl K3.toFinFlag FlagAlgebra_4_0_0_8
  apply forbidLE_trans quadraticForm_trans
  apply forbidLE_trans_forbidEq_right ?_
        (forbidEq_smul (forbidEq_symm
          (one_forbidEq_forbidExpand_one K3.toFinFlag 4)))
  simp [flagQuadraticForm, v₁, M₁_real, ratMatrixToReal, M₁,
        Fin.sum_univ_four, add_assoc]
  simp [v₂, M₂_real, ratMatrixToReal, M₂, Fin.sum_univ_three]
  reduce_downward_flagmul
  expand_one_at 4
  simp [smul_smul, downward_add, downward_smul]
  ac_sort_rhs_pipeline
  apply forbidLE_of_le
  flag_nonneg
\end{lstlisting}

Each line corresponds directly to one step of the flag algebra method.
The \texttt{have}-block adds each block's quadratic form to the
left-hand side under the forbid relation; the
\texttt{apply forbidLE\_trans\_forbidEq\_right} step replaces the unit
$1$ by its size-$N$ expansion as a forbid-equivalence
(\texttt{expand\_one\_at $N$} fills it in later); the two
\texttt{simp} lines unfold the quadratic forms into explicit sums of
products of $\sigma$-flags; \texttt{reduce\_downward\_flagmul}
(Section~\ref{sec:tactics}) substitutes the precomputed
\texttt{flagMul\_*} expansion for every product; \texttt{expand\_one\_at
$N$} finishes the right-hand-side rewrite; the closing pipeline
normalizes scalars and sorts; and \texttt{forbidLE\_of\_le} together
with \texttt{flag\_nonneg} discharges the resulting non-negativity
inequality.

\paragraph{What is parametrized.}
The translation script chooses the matching
\texttt{Fin.sum\_univ\_<n>} for each block's dimension (\texttt{two},
\texttt{three}, \ldots, \texttt{eight}), the argument $N$ to
\texttt{expand\_one\_at}, the size argument to
\texttt{one\_forbidEq\_forbidExpand\_one}, and the names $M_t$, $v_t$ in
each \texttt{simp}. The order of \texttt{apply forbidLE\_add\_-
QuadraticForm} calls inside the \texttt{have} is reversed relative to
the addition order in the \texttt{have} statement, so the outermost
$+$ peels off first; the script emits them in reverse without further
configuration.

\subsection{Branch B: $n_{\mathit{obj}} < N$}

When the objective lives below the host size -- as in Mantel
($n_{\mathit{obj}} = 2 < 3 = N$) or K4turan
($n_{\mathit{obj}} = 2 < 4 = N$) -- the body needs an extra
\emph{forbid-conditioned expansion} step that lifts the objective up to
host size, dropping forbidden contributions in the process. The
translation script emits an auxiliary lemma whose name is
\texttt{<theorem>\_expand\_under\_forbid}, and inserts a single
\texttt{rw} after the unit-expansion step:
\begin{lstlisting}
apply forbidLE_trans_forbidEq_right ?_
      (forbidEq_smul (forbidEq_symm
        (one_forbidEq_forbidExpand_one K3.toFinFlag 3)))
rw [forbidLE_rw_left_add_right mantel_flagAlgebra_expand_under_forbid]
\end{lstlisting}
The auxiliary lemma, for Mantel, is

\begin{lstlisting}
lemma mantel_flagAlgebra_expand_under_forbid
    : FlagAlgebra_2_0_0_1 =[K3.toFinFlag]
        (1 / 3 : ℝ) • FlagAlgebra_3_0_0_1
      + (2 / 3 : ℝ) • FlagAlgebra_3_0_0_2 := by
  have h_unit_3 : (FlagAlgebra_3_0_0_3 : FlagAlgebra ∅ₜ)
      = ⟦unitVector (⟨3, Flag_3_0_0_3⟩ : FinFlag ∅ₜ)⟧
    := (Quotient.out_inj.mp rfl).symm
  have h_zero_3 : (FlagAlgebra_3_0_0_3 : FlagAlgebra ∅ₜ) =[K3.toFinFlag] 0 := by
    rw [h_unit_3]
    apply unitVector_forbidEq_zero
    rw [unlabel_emptyType]
    exact lt_of_le_of_ne
      (flagListDensity₁_ge_zero K3.toFinFlag.2 Flag_3_0_0_3)
      (Ne.symm flagDensity1_K3_Flag_3_0_0_3_ne_zero)
  have h_eq : FlagAlgebra_2_0_0_1 =[K3.toFinFlag]
      (1 / 3 : ℝ) • FlagAlgebra_3_0_0_1
    + (2 / 3 : ℝ) • FlagAlgebra_3_0_0_2
    + FlagAlgebra_3_0_0_3 :=
    forbidEq_of_eq (by prove_flag_expand 3)
  rw [forbidEq_rw_right_add_left h_zero_3, add_zero] at h_eq
  exact h_eq
\end{lstlisting}

The structure is:
\begin{enumerate}
  \item For every \emph{forbidden} host flag with non-zero objective
    density, identify the flag with its unit-vector representative
    (\texttt{h\_unit\_i}) and prove that it is forbid-equivalent to
    zero (\texttt{h\_zero\_i}) using the
    \texttt{flagDensity1\_<H>\_..\_ne\_zero} lemma registered in
    Part~2.
  \item Prove the full unconditioned expansion
    $F_{\mathit{obj}} = \sum_j d_j\,F_j^{(N)}$ -- including the
    forbidden hosts -- with a single call to
    \texttt{prove\_flag\_expand $N$}. The list of coefficients $d_j$
    is computed by the script as the induced density of $F_{\mathit{obj}}$
    in each host graph (a brute-force enumeration over
    $\binom{N}{n_{\mathit{obj}}}$ subsets) and emitted as part of the
    lemma statement.
  \item Drop the forbidden summands by rewriting each
    \texttt{h\_zero\_i} into the unconditioned expansion lifted to
    forbid-equivalence, followed by \texttt{add\_zero}.
\end{enumerate}

\paragraph{Density evaluation table.}
The closing step of \texttt{prove\_flag\_expand $N$} needs every
\texttt{flagDensity}$_1\,\texttt{Flag\_obj}\,\texttt{Flag\_host\_i}$ to
evaluate to a concrete rational by \texttt{simp}. The translation script
emits the full table of these as \texttt{@[simp]} lemmas above the
auxiliary lemma, each closed by \texttt{native\_decide} after a small
\texttt{dsimp} / \texttt{rw [flagDensity\textsubscript{1}\_eq\_sym2EmptyTypeFlagDensity\textsubscript{1}]} normalization.
The coefficients in the lemma statement are read out of the same
induced-density computation, so the table values and the expansion
coefficients are consistent by construction. In Mantel, for instance,
the script emits four such lemmas (one per size-$3$ host),
\texttt{auto\_flagDensity1\_2\_0\_0\_1\_3\_0\_0\_\{0,1,2,3\}}, evaluating
to $0$, $1/3$, $2/3$, and $1$ respectively.

\paragraph{Right-associativity for multiple blocks.}
The rewrite \texttt{forbidLE\_rw\_left\_add\_right} only matches the
pattern $\mathit{obj} + \texttt{?}$ at the top-level $+$. When there
are two or more blocks the left-hand side arrives left-associated as
$((\mathit{obj} + Q_1) + Q_2) + \ldots$, so the pattern does not fire.
The translation script inserts a \texttt{simp only [add\_assoc]} before
the rewrite to right-associate the sum first; for the single-block case
this would be a no-op, so the script emits it only when $T \ge 2$.

\section{Trust model}
\label{sec:trust}

The translation script computes a substantial amount: it isomorphism-
matches flagmatic strings to canonical indices, assembles $R_t \cdot
Q_t' \cdot R_t^\top$ in exact rationals, factors that as $L \cdot
\mathrm{diag}(d) \cdot L^\top$, exhaustively counts induced subgraphs to
fill the branch-B density table, and chooses which Lean
identifiers and which \texttt{load\_*} files to emit. None of these
computations enters the trusted base.

Every emitted lemma is closed by one of:
\begin{itemize}
  \item \texttt{decide +kernel} or \texttt{decide}, which re-evaluates
    the equality inside the Lean kernel,
  \item \texttt{native\_decide}, which compiles the decidable statement
    to native code and re-evaluates it (still inside the kernel for
    soundness purposes), or
  \item one of the verified tactics of Section~\ref{sec:tactics}
    (\texttt{prove\_flag\_expand}, \texttt{reduce\_downward\_flagmul},
    \ldots), each of which uses only the \texttt{flagSet}, \texttt{Flag},
    and \texttt{flagMul} names registered by the meta-commands of
    Section~\ref{sec:density-tables}.
\end{itemize}
A miscomputation in the script -- a wrong LDL pair, a wrong density, a
mis-resolved identifier -- surfaces as a compile error, never as a
silently accepted false claim. The trust model here is exactly the one
already in place for the meta-commands and the data-generation
scripts of Section~\ref{sec:flag-defs}, and the translation script
fits into that model without adding to it.

\paragraph{The dependency check as a safety net.}
The strict \texttt{check-deps} behavior described in Section~7 is a
defence-in-depth measure for one specific failure mode that the
re-verification does not by itself prevent: silently loading the
\emph{wrong} density JSON. If two density files have the same
``host\_from\_pattern'' base name but one was computed under the
$K_3$-forbid and the other under the $K_5$-forbid, then loading the wrong
one yields fully verified theorems that are simply theorems about a
different problem. Insisting on the explicit \texttt{\_forbid\_<H>}
suffix makes this confusion impossible at the file-system level.

\section{Verified scenarios}
\label{sec:scenarios}

Six certificates have been put through the pipeline and produce files
that the Lean kernel accepts with no \texttt{sorry}.

\begin{center}
\begin{tabular}{lcccccc}
\toprule
Certificate & Forbid & $N$ & $n_{\mathit{obj}}$ & Branch & $T$ & Bound \\
\midrule
Mantel               & $K_3$ & 3 & 2 & B & 1 & $1/2$    \\
K3forbidP3           & $K_3$ & 3 & 3 & A & 1 & $3/4$    \\
K3forbidC4           & $K_3$ & 4 & 4 & A & 2 & $3/8$    \\
K4turan              & $K_4$ & 4 & 2 & B & 2 & $2/3$    \\
ErdosPentagon        & $K_3$ & 5 & 5 & A & 3 & $24/625$ \\
K5turan              & $K_5$ & 5 & 2 & B & 4 & $3/4$    \\
\bottomrule
\end{tabular}
\end{center}

The six cases between them exercise every feature of the translation
script.
\begin{itemize}
  \item \textbf{Mantel} ($K_3$, $N{=}3$, branch B, single block): the
    smallest non-trivial case, used to validate the branch-B template
    (auxiliary lemma, induced-density table, forbidden-term dropping)
    in isolation.
  \item \textbf{K3forbidP3} ($K_3$, $N{=}3$, branch A, single block):
    the branch-A counterpart of Mantel at the same host size; the
    objective is the size-$3$ path and lives at host size, so no
    expansion lemma is needed.
  \item \textbf{K3forbidC4} ($K_3$, $N{=}4$, branch A, two blocks):
    exercises branch A's per-block structure and the order in which the
    quadratic-form additions are emitted.
  \item \textbf{K4turan} ($K_4$, $N{=}4$, branch B, two blocks): the
    first instance with $H \ne K_3$; verifies that the
    \texttt{CommonGraphs.lean} parser-driven forbid resolution
    (Section~\ref{sec:mapping}) succeeds and that
    \texttt{load\_forbid\_density\_theorems} and the proof body work
    against a non-$K_3$ forbid.
  \item \textbf{ErdosPentagon} ($K_3$, $N{=}5$, branch A, three blocks):
    the largest branch-A case in the suite, with three quadratic forms
    of dimensions $8$, $6$, $5$ -- a single
    \texttt{reduce\_downward\_flagmul} call dispatches all
    $8^2 + 6^2 + 5^2 = 125$ product terms.
  \item \textbf{K5turan} ($K_5$, $N{=}5$, branch B, four blocks): the
    first instance with $H = K_5$, and the largest case overall. Four
    quadratic forms of dimension $8$ each yield $4 \times 8^2 = 256$
    product terms, all dispatched by a single
    \texttt{reduce\_downward\_flagmul}, and the branch-B auxiliary
    lemma carries an induced-density table of size $|\text{graphs}_5|
    \times 1$.
\end{itemize}

\section{Adding a new forbidden graph}
\label{sec:extension}

The translation script \texttt{flagmatic\_to\_lean.py} itself does not
need to be changed to support a new forbidden graph $X$: it discovers
which forbid tags exist by re-parsing \texttt{CommonGraphs.lean} on
every run, so adding an entry there is enough to make $X$ visible to the
identifier-mapping pass of Section~\ref{sec:mapping}. The Lean-side
meta-commands that turn the JSON data into lemmas, on the other hand,
still currently dispatch on the forbid tag explicitly (one branch per
$K_3$, $K_4$, $K_5$), so they need a new branch for $X$. The full
extension procedure is:
\begin{enumerate}
  \item In \texttt{Forbid/CommonGraphs.lean} add the two lines
\begin{lstlisting}
def X : SimpleGraph (Fin n) := <X's definition>
lemma X_toFinFlag_eq : X.toFinFlag = ⟨n, Flag_n_0_0_i⟩ := by <...>
\end{lstlisting}
    where $i$ is the canonical index of $X$ inside
    \texttt{graphs\_n.json}. The patterns for $K_3$, $K_4$, $K_5$ in
    that file give templates to follow.
  \item Generate the forbid-conditioned data files by running
\begin{lstlisting}[language={}]
python gen_free_indices.py --forbid-X
python calculate_densities.py --host <hostFlags>.json \
       --pattern <patternFlags>.json --forbid-X
\end{lstlisting}
    once for each $(\text{host},\text{pattern})$ pair the new theorem
    will need. The output names automatically carry the
    \texttt{\_forbid\_X} suffix.
  \item Add a branch for $X$ to the meta-commands in
    \texttt{Flags/Densities/DensityLoader.lean} (and, if the new
    theorem uses products, to
    \texttt{Flags/Densities/MulLoader.lean}). The existing
    \texttt{K3}/\texttt{K4}/\texttt{K5} branches give templates that
    can be copied verbatim and renamed; each branch is roughly $15$
    lines, and the only $X$-specific content is the name $X$ itself
    and the corresponding flag index $i$ from Step~1. When we added
    $K_5$ support during this work, the changes to
    \texttt{DensityLoader.lean} were exactly of this shape -- a single
    new branch parallel to the existing ones.
  \item Run \texttt{gen-skeleton} on the new certificate, then
    \texttt{lake build}.
\end{enumerate}

Why the meta-commands need this branch is worth a brief comment. Each
\texttt{flagDensity}$_1$ lemma the meta-command emits has its forbid
graph baked into the proof script (\texttt{rw [X.toFinFlag\_eq]; \ldots
native\_decide}), so a generic version would need to abstract over the
\texttt{toFinFlag\_eq} family. A future cleanup could do this; at the
present scale of K3, K4, K5 the manual copy-paste cost has not been
high enough to motivate the abstraction.

The K5turan certificate has been carried through all four steps and
appears in Section~\ref{sec:scenarios} as the kernel-verified end-to-end
example of this extension procedure: the $K_5$ branches in
\texttt{CommonGraphs.lean} and \texttt{DensityLoader.lean} are in place,
the K5-specific JSON files have been produced, and the generated
\texttt{K5turan.lean} compiles with no \texttt{sorry}.

\section{What is still not automated}

The translation pipeline closes the gap between a flagmatic
certificate and a kernel-verified Lean proof, but two caveats are worth
stating about its current coverage.

\paragraph{Branch-B pathological cases.}
The induced-density table that branch B emits is closed by
\texttt{native\_decide}, and the auxiliary lemma's final step is
\texttt{prove\_flag\_expand $N$}. Both work on every example we have
encountered, but \texttt{prove\_flag\_expand} ends in a
\texttt{ring\_nf} step, and we have not formally characterized the
family of expansions it is guaranteed to close. A pathological
certificate -- one whose induced-density coefficients, after lifting
through \texttt{forbidEq\_of\_eq}, produce a form
\texttt{ring\_nf} fails to normalize -- could in principle require
manual intervention here. We have not run into such a case among the
six verified scenarios of Section~\ref{sec:scenarios}, but we have
also not proven that none exists.

\paragraph{Non-$K_n$ forbidden graphs.}
The identifier-mapping layer of Section~\ref{sec:mapping} is general:
it iso-matches the certificate's forbid string against every entry of
\texttt{CommonGraphs.lean} via vertex permutation, with no assumption
that the forbid graph is a complete graph. The Lean side is general in
the same sense: nothing in the proof body or the load commands assumes
$X = K_n$. In practice, however, only $K_3$, $K_4$, $K_5$ are currently
defined in \texttt{CommonGraphs.lean} and matched by branches in the
meta-commands, so a forbid such as $C_4$ or $P_4$ has not yet been
exercised end to end. Section~\ref{sec:extension} describes the
extension procedure; once a non-$K_n$ case is carried through it will
also serve as the first empirical evidence that the
\emph{generality} we claim here is in fact realized.

\section{Summary}

The framework presented here takes a flagmatic SDP certificate, in the
form of a single JSON file, and produces a Lean module that the kernel
accepts as a proof of the corresponding density inequality. The
intermediate steps -- the parsing of flagmatic strings into canonical
indices, the assembly and $\mathrm{LDL}^\top$ factorization of each
block, the choice of forbid-tagged density files, the unfolding of the
unit element under the forbid relation, and the dropping of
forbidden host flags in the branch-B auxiliary lemma -- are all
performed by the translation script, but none of them are admitted into
the trusted base: every numerical fact is re-checked by
\texttt{decide}/\texttt{native\_decide} or discharged by a verified
tactic, and a mismatch between script and kernel is reported as a
compile error rather than a silently accepted theorem.

The framework rests on three layers that were built together. The
generated flag definitions of Section~\ref{sec:flag-defs} fix a single
naming convention; the density tables of
Section~\ref{sec:density-tables} pour numerical facts into that
convention; the tactics of Section~\ref{sec:tactics} consume them. On
top of these the translation script of
Sections~\ref{sec:mapping}--\ref{sec:body} performs the
certificate-to-Lean translation. Extending the system to a new
forbidden graph reduces to two lines in \texttt{CommonGraphs.lean} and
a script run; extending it to a new theorem with an existing forbid
reduces to one command.

The six certificates of Section~\ref{sec:scenarios} -- Mantel,
K3forbidP3, K3forbidC4, K4turan, Erd\H{o}s pentagon, and K5turan --
between them exercise every feature of the pipeline, spanning host sizes
$N \in \{3, 4, 5\}$, block counts $1$ through $4$, both branches, and
forbids $K_3$, $K_4$, $K_5$. None of these required hand transcription
of the certificate into Lean.

\end{document}
